\section{The odd-degree rank bound and the finite-image estimate}
\label{sb-section}

The remaining homology comparison has a short proof using the
classical rational homology of the algebraic \(K\)-theory space.
We state those classical inputs explicitly.  All subsequent
integral counting arguments are proved below, with explicit
constants.  The method uses algebraic \(K\)-theory in place of
the differential-form argument of~\cite{alpo-uniform-source}.

Let \(L\) be a totally imaginary number field, put
\(R=\mathcal O_L\), and write \(r=r_2(L)\).
For \(N,j\ge0\), retain the notation
\[
 t_L(N,j)=\#H_j(\operatorname{GL}_N(R),\mathbb Z)_{\rm tors}.
\]
The finiteness needed here was proved in
Sections~\ref{ac-section} and~\ref{ht-section}.
All homology has trivial coefficients unless otherwise stated.

\subsection{Stable homology and the repeated direct-sum map}

As in Section~\ref{kt-section}, put
\(X=B\operatorname{GL}(R)^+\), the connected component of the
algebraic \(K\)-theory infinite loop space.  Its multiplication
is induced by direct sum of matrices.  It is homotopy
associative and homotopy commutative, has a unit, and
\[
 H_j(X,\mathbb Z)=H_j(\operatorname{GL}(R),\mathbb Z),
 \qquad \pi_j(X)=K_j(R)\quad(j\ge1).
 \tag{SB1}
\]
We use the following established results.

\begin{enumerate}
\item Integral homological stability identifies
\(H_j(\operatorname{GL}_N(R),\mathbb Z)\) with
\(H_j(X,\mathbb Z)\) for
\(N\ge\max(2j+1,j+2)\).
The stable-range condition for \(R\) was proved in
Section~\ref{kt-section}.
\cite[IV, Corollary 1.14.1]{weibel-kbook}
\item For a path-connected, homotopy associative \(H\)-space,
its rational homology is generated, under its homology
product, by primitive elements.  The primitive elements
in positive degree are precisely the image of its rational
homotopy under the Hurewicz map, and that map is injective.
Here an element \(v\) is primitive when
\(\Delta v=v\otimes1+1\otimes v\), where \(\Delta\)
comes from the diagonal of the space.
\cite[Corollary 4.18, Theorem 5.18, and Appendix]{milnor-moore}
\item For every odd \(j\ge3\),
\(\dim_{\mathbb Q}(K_j(R)\otimes\mathbb Q)=r\).
For \(j=1\), this dimension is \(r-1\): the classical
identity \(K_1(R)=R^\times\) and
Theorem~\ref{su-unit-theorem} give this assertion.
The identity in degree one is the Bass--Milnor--Serre
theorem, not an assertion about arbitrary Dedekind domains.
\cite[III, Example 1.3.6; IV, Theorems 1.17--1.18]{weibel-kbook}
\end{enumerate}

These are classical inputs with their stated hypotheses;
we do not reprove the general stability or rational
homotopy theorems here.  In particular, no quantitative
comparison of integral \(K\)-groups with homology is being
assumed in this section.

For a positive integer \(d\), let
\[
 \psi_{d,N}:\operatorname{GL}_N(R)
       \longrightarrow\operatorname{GL}_{dN}(R),\qquad
       g\longmapsto\operatorname{diag}(g,\ldots,g).
\]
Its stable homology map will be denoted by \(\Psi_d\).
The choice of an ordering of the blocks changes the map
by conjugation and therefore does not change homology.
On \(X\), this is the map which adds a point to itself
\(d\) times under the direct-sum multiplication.

\begin{lemma}\label{sb-weight}
Let \(m\ge1\).  In rational stable homology, the image of
\[
       \prod_{j=2}^{m}(\Psi_d-d^j)
           :H_m(X,\mathbb Q)\longrightarrow H_m(X,\mathbb Q)
 \tag{SB2}
\]
is contained in the space of primitive elements of degree
\(m\).  An empty product, when \(m=1\), means the identity.
\end{lemma}
\begin{proof}
If \(v\) is primitive, its iterated diagonal is the sum
of the \(d\) tensors having \(v\) in one position and
\(1\) elsewhere.  Multiplication therefore gives
\(\Psi_d(v)=dv\).  Homotopy commutativity implies that
the repeated-sum map preserves the homology product:
equivalently, the product and the diagonal in the graded
commutative homology algebra show that
\(\Psi_d(vw)=\Psi_d(v)\Psi_d(w)\).
Consequently a product of \(s\) positive-degree
primitive elements has eigenvalue \(d^s\).

Primitive generation expresses every element in degree
\(m\) as a sum of such products.  Every factor has
degree at least one, so \(1\le s\le m\).
If \(s\ge2\), the factor \(\Psi_d-d^s\) in (SB2)
annihilates the product.  The remaining terms are single
primitive elements of degree \(m\).  This also proves
the statement for \(d=1\); no division by an eigenvalue
difference has been used.
\end{proof}

We now keep track of the finite matrix sizes.  Fix
\(d\ge1\), \(N\ge1\), and an integral homology class
\(x\in H_m(\operatorname{GL}_N(R),\mathbb Z)\).
Let \(N_t=d^tN\), let \(s_t\) be the usual inclusion
from size \(N_t\) to size \(N_{t+1}\), and define
\[
 \begin{split}
 Z_{d,0}(x)&=x,\\
 Z_{d,t+1}(x)&=(\psi_{d,N_t})_*Z_{d,t}(x)
                         -d^{t+2}(s_t)_*Z_{d,t}(x).
 \end{split}
 \tag{SB3}
\]
These are homomorphisms of integral homology groups.
They are exactly the repeated direct-sum differences
used by the source.

\begin{proposition}\label{sb-odd-rank}
If \(m\) is positive and odd and \(N\ge2m+4\), then
\[
 \operatorname{rank}_{\mathbb Z}
 \operatorname{im}\bigl(Z_{d,m-1}:
 H_m(\operatorname{GL}_N(R),\mathbb Z)
 \longrightarrow H_m(\operatorname{GL}_{d^{m-1}N}(R),\mathbb Z)
                         \bigr)\le r .
 \tag{SB4}
\]
The extra hypothesis \(m<r\) in the submitted assertion
is not needed for this rank estimate.
\end{proposition}
\begin{proof}
All sizes in (SB3) are at least \(N\), since \(d\ge1\).
They lie in the integral stable range.  After these
identifications and tensoring with \(\mathbb Q\),
\(Z_{d,m-1}\) becomes (SB2).
Its image is primitive by Lemma~\ref{sb-weight}.
The dimension of that space is \(r\) for odd
\(m\ge3\), and \(r-1\) for \(m=1\), by the classical
inputs above.  Tensoring with \(\mathbb Q\) is exact,
so it computes the rank of the integral image.
This proves (SB4).
\end{proof}

\subsection{Two elementary counting lemmas}

\begin{lemma}\label{sb-finite-rank}
Let \(H\) be an abelian group with finite torsion subgroup,
let \(W\subset H\) have rank at most \(r\), and let
\(f:H\to B\) be a homomorphism to an abelian group
annihilated by a positive integer \(q\).  Every finite
subset of \(f(W)\) has cardinality at most
\[
                     q^r\,\#H_{\rm tors}.
 \tag{SB5}
\]
In particular, \(f(W)\) is finite with that bound.
Finite generation of \(H\) is not required.
\end{lemma}
\begin{proof}
Lift a given finite subset to finitely many elements of
\(W\), and let \(W'\) be the subgroup they generate.
The structure theorem for finitely generated abelian
groups gives
\(W'\simeq\mathbb Z^s\oplus T'\), with
\(s\le r\) and \(T'\subset H_{\rm tors}\).
The image of the free part has at most \(q^s\) elements,
and the image of \(T'\) has at most \(\#H_{\rm tors}\)
elements.  This proves (SB5).
If \(f(W)\) had more elements, choosing a finite subset
of one greater than the bound would contradict (SB5).
\end{proof}

\begin{lemma}\label{sb-coefficients}
For every group \(G\), positive integer \(q\), and
\(m\ge1\), there is a natural exact sequence
\[
 0\longrightarrow H_m(G,\mathbb Z)/qH_m(G,\mathbb Z)
 \longrightarrow H_m(G,\mathbb Z/q)
 \xrightarrow{\ \beta\ }H_{m-1}(G,\mathbb Z)[q]
 \longrightarrow0.
 \tag{SB6}
\]
Here \(H[q]=\{h\in H:qh=0\}\).
\end{lemma}
\begin{proof}
Use the integral bar complex \(C_\bullet\), whose groups
are free abelian, and its quotient \(C_\bullet/qC_\bullet\).
Lift a cycle \(\overline z\) of degree \(m\) to
\(z\in C_m\).  There is \(w\in C_{m-1}\) with
\(\partial z=qw\).  Since \(C_{m-2}\) has no
\(q\)-torsion, \(\partial w=0\); for \(m=1\)
use the ordinary complex with \(\partial:C_0\to0\).
Set \(\beta[\overline z]=[w]\).
Changing the lift by \(qu\) changes \(w\) by
\(\partial u\), and changing the cycle by a boundary
has the same harmless effect.  Thus \(\beta\) is a
well-defined homomorphism.  Its value is killed by
\(q\), because \(qw=\partial z\).

If \([w]=0\), write \(w=\partial u\).  Then
\(z-qu\) is an integral cycle reducing to
\(\overline z\), proving exactness in the middle.
Conversely every integral cycle has zero image under
\(\beta\).  If an integral cycle \(z\) reduces to a
boundary, write \(z=\partial v+qu\).  Applying
\(\partial\) and using the freeness of the chain groups
gives \(\partial u=0\).  Therefore the kernel of
reduction on integral homology is exactly \(qH_m\).
Finally, if \([w]\in H_{m-1}[q]\), choose a cycle
\(w\) and a chain \(z\) with \(qw=\partial z\).
The reduction of \(z\) is a cycle mapping to \([w]\).
This proves surjectivity and every term of (SB6).
\end{proof}

\subsection{The finite-image estimate}

Fix a prime \(\ell\), an integer \(\nu\ge0\), and
\(A=\mathbb Z/\ell^\nu\mathbb Z\).  Let \(B\) be
an \(A\)-module, and suppose for each \(N\ge0\) that
\[
 \Phi_N:C_m(\operatorname{GL}_N(R);A)\longrightarrow B
\]
is \(A\)-linear and vanishes on boundaries.
Assume that, on cycles, it is unchanged by stabilization
and satisfies
\[
     \Phi_{dN}((\psi_{d,N})_*y)=d\Phi_N(y)
                         \qquad(d\ge0).
 \tag{SB7}
\]
These are the precise properties required of the
family; they are proved for the finite-field functionals
in Section~\ref{sf-section}.

\begin{theorem}\label{sb-upper}
Let \(m\ge1\) be odd and \(N_0=2m+4\).
Every finite subset \(S\subset B\) whose elements have
the form
\[
                \ell^{m-1}\Phi_N(y),
       \qquad \partial y=0,
 \tag{SB8}
\]
where \(N\) and \(y\) may depend on the element,
satisfies
\[
 \#S\le (\ell^\nu)^r\,
       t_L(N_0,m)\,t_L(N_0,m-1).
 \tag{SB9}
\]
As in Proposition~\ref{sb-odd-rank}, the extra
source hypothesis \(m<r\) can be omitted here.
\end{theorem}
\begin{proof}
If \(\nu=0\), then \(A\) is the zero ring and every
unital \(A\)-module is zero.  The assertion follows.
Suppose henceforth \(\nu\ge1\), and write
\(q=\ell^\nu\).
There are finitely many elements of \(S\), so choose
one size \(N\ge N_0\) larger than all their given sizes.
Stabilization expresses them all in the form (SB8)
at this common size.

For each size, reduction of integral cycles and then
\(\Phi\) defines a homomorphism
\[
 \phi_N:H_m(\operatorname{GL}_N(R),\mathbb Z)
                        \longrightarrow B.
\]
It is well-defined because integral boundaries reduce
to boundaries.  It inherits stabilization and (SB7).
Applying it to (SB3), with \(d=\ell\), gives
\[
 \phi_{\ell^{m-1}N}(Z_{\ell,m-1}(x))
    =\prod_{j=2}^{m}(\ell-\ell^j)\phi_N(x)
    =\ell^{m-1}u\phi_N(x),
 \quad
 u=\prod_{j=2}^{m}(1-\ell^{j-1}).
 \tag{SB10}
\]
Each factor of \(u\) is congruent to \(1\) modulo
\(\ell\), so \(u\) is invertible modulo \(q\).
Choose \(v\in\mathbb Z\) with \(uv\equiv1\pmod q\).
It follows from (SB10) that every
\(\ell^{m-1}\phi_N(x)\) belongs to
\[
 \phi_{\ell^{m-1}N}
       (\operatorname{im}Z_{\ell,m-1}),
\]
by replacing \(x\) in the right-hand side by \(vx\).
Proposition~\ref{sb-odd-rank} bounds the rank of this
image by \(r\).  Integral stability identifies the
torsion subgroup at size \(\ell^{m-1}N\) with the one
at size \(N_0\).  Lemma~\ref{sb-finite-rank} therefore
bounds the image of all integral cycles, after
multiplication by \(\ell^{m-1}\), by
\(q^r t_L(N_0,m)\).

It remains to account for cycles modulo \(q\) which
need not lift to integral cycles.  Apply (SB6) to
\(G=\operatorname{GL}_N(R)\).
Its right-hand group is a subgroup of the integral
torsion in degree \(m-1\), of order at most
\(t_L(N_0,m-1)\) by stability.
The homomorphism induced by
\(\ell^{m-1}\Phi_N\) maps each coset of the kernel
of \(\beta\) to a translate of the image on that
kernel.  The preceding paragraph bounds that image
by \(q^r t_L(N_0,m)\), and there are at most
\(t_L(N_0,m-1)\) such cosets.  Multiplication of the
two bounds gives (SB9).
\end{proof}

The factor involving degree \(m-1\) is necessary in
this argument: it measures the failure of a cycle
modulo \(q\) to lift to an integral cycle.
Likewise, the repeated direct-sum differences in
(SB3) are essential to isolate the primitive part.
For example, take \(L=\mathbb Q(\zeta_{13})\), so
\(r=6\).  Its stable rational homology in degree three
has dimension \(6+\binom{5}{3}=16\), whereas (SB4)
bounds the rank after the two differences by \(6\).
Indeed the primitive degree-one space has dimension
\(5\), the degree-three space has dimension \(6\),
and the rational homology algebra is the graded
symmetric algebra on its primitives by the same
classical Hopf-algebra theorem.  There are no
degree-two primitives, by Borel's rank theorem.
Thus the extra ten classes are the triple products
in degree one; the factor \(\Psi_d-d^3\) removes them.
This example refutes the unfiltered rank statement,
not the statement used by the submitted proof.

The arguments above prove Proposition~\ref{sb-odd-rank} and
Theorem~\ref{sb-upper} using the explicitly stated classical inputs.


